\chapter{Investment}\label{ch:b2:investment}

% figdata: recession (inv -31.4, cons -1.9); iy_summary (USA late 20.3)
Investment is the most volatile part of spending. In the United States it has been about a fifth of GDP on average since
1950, and from the end of 2007 to the middle of 2009 real private investment fell by 31 per cent while real consumption fell
by 2 per cent. The theory of this chapter says that a firm invests when the market values its capital above
what the capital costs to replace. The chart puts that to a first test on the data Grunfeld collected on eleven large
American firms from 1935 to 1954.

\begin{openingchart}{Gross investment against the market value of the firm, eleven United States firms, 1935--1954, each
  divided by the firm's own average over the period (220 firm-years).}{statsmodels bundled data (public domain), grunfeld
  (Grunfeld's data as recreated by Kleiber and Zeileis); our computation}
\begin{tikzpicture}
\begin{axis}[omrecord, width=0.62\linewidth, height=0.5\linewidth, xmin=0, xmax=2.5, ymin=0, ymax=2.6,
    xlabel={Market value, relative to the firm's average}, ylabel={Investment, relative to the firm's average}]
  \addplot[only marks, mark=*, mark size=1.2pt, omDef, fill opacity=0.6]
    table[x=v_rel, y=i_rel] {figdata/bachelor-2/20-investment/out/grunfeld.csv};
  \addplot[omThm, thick, domain=0:2.5, samples=2] {0.93*x + 0.07};
\end{axis}
\end{tikzpicture}
\end{openingchart}

% figdata: grunfeld_summary (n 220, firms 11, corr 0.57, slope 0.93)
In years in which a firm's market value was high relative to its own average, it invested more: the correlation over the
220 firm-years is 0.57, and the fitted line has a slope of 0.93. But the cloud is wide: at the same relative value, some
firms invested half their average and others twice it. This chapter builds the theory that predicts the slope, and
asks what else drives investment: uncertainty, which makes waiting valuable, and finance, whose role is disputed.

\begin{recall}
From the Year 1 volume: the rental price of capital and the marginal product of capital. From \cref{ch:b2:asset-prices}: the
value of an asset as the present value of what it pays. From the mathematics volumes: systems of two differential
equations and their phase diagrams.
\end{recall}

\begin{notation}[Symbols of this chapter]\label{not:b2:investment:symbols}
$K$ is the capital stock and $I$ gross investment; $\delta$ the depreciation rate; $r$ the real interest rate; $p_K$ the price of a
unit of capital; $q$ the value of a unit of installed capital relative to its price; $\chi$ the adjustment-cost parameter.
\end{notation}

\section{The user cost of capital}\label{sec:b2:investment:user-cost}

A firm that buys a machine at price $p_K$ and holds it for a year gives up interest $rp_K$, loses $\delta p_K$ to wear, and gains $\pi_K p_K$ if the
price of machines rises at rate $\pi_K$.

\begin{definition}[User cost of capital]\label{def:b2:investment:user-cost}
The \emph{user cost of capital}\index{user cost of capital} is the cost of holding one unit of capital for one period:
$p_K(r + \delta - \pi_K)$, where $\pi_K$ is the rate of change of the price of capital goods.
\end{definition}

\begin{definition}[Desired capital stock]\label{def:b2:investment:desired-capital}
The \emph{desired capital stock}\index{desired capital stock} is the stock at which the value of the marginal product of
capital equals its user cost.
\end{definition}

With output $Y = AK^\alpha L^{1-\alpha}$ and output price 1, the marginal product is $\alpha Y/K$, so the desired stock is $K^* = \alpha Y/(p_K(r + \delta - \pi_K))$. With $\alpha = 0.3$, $Y = 1{,}000$, $p_K = 1$,
$r = 0.05$, $\delta = 0.10$ and $\pi_K = 0$, the user cost is 0.15 and $K^* = 2{,}000$. A rise in the interest rate to 0.07 raises the user cost to 0.17 and lowers $K^*$ to 1,765.
Taxes enter the same way: a tax on profits lowers the after-tax value of the marginal product, and allowances for
depreciation or credits for investment lower the effective price $p_K$.

The \omterm{def:b2:investment:desired-capital}{desired capital stock} says where the firm wants to be, not how fast it gets there. If the firm could adjust at no
cost, a fall in the interest rate would produce an instant jump in the capital stock, an infinite rate of investment for
a moment. Investment is in fact spread over years.

\section{Adjustment costs and Tobin's \texorpdfstring{$q$}{q}}\label{sec:b2:investment:q}

\begin{definition}[Adjustment cost]\label{def:b2:investment:adjustment-cost}
An \emph{adjustment cost}\index{adjustment cost} is a cost of changing the capital stock over and above the price of the
capital goods, such as installation, disruption and training, which rises more than in proportion to the speed of
investment.
\end{definition}

Let the cost of investing $I$ with capital $K$ be $I + \tfrac{\chi}{2}(I/K - \delta)^2K$: the price of the goods, plus a convex cost that is zero when the firm only replaces
worn capital. Let $q$ be the value to the firm, in present value, of one more unit of installed capital, measured relative to
its price $p_K = 1$. The firm invests until the marginal cost of a unit, $1 + \chi(I/K - \delta)$, equals $q$:
\begin{equation}\label{eq:b2:investment:q}
\frac{I}{K} = \delta + \frac{q - 1}{\chi}.
\end{equation}
With $\chi = 10$ and $\delta = 0.08$, a firm with $q = 1.3$ invests 11 per cent of its capital a year; at $q = 1$ it only replaces what wears out.

\begin{definition}[Tobin's q]\label{def:b2:investment:tobins-q}
\emph{Tobin's q}\index{Tobin's q} is the ratio of the value of capital to its replacement cost. \emph{Marginal $q$} is the value
of one more unit of installed capital over its price; \emph{average $q$} is the market value of the firm over the
replacement cost of its whole capital stock.
\end{definition}

Marginal $q$ drives investment but cannot be observed; average $q$ can, from stock prices. When the firm has constant returns to
scale in production and in \omterm{def:b2:investment:adjustment-cost}{adjustment costs} and takes prices as given, the two are equal, a result this book
admits. The theory then says: regress investment on average $q$ and nothing else should matter.

\begin{model}{$q$ theory of investment}
\assumes Firms that maximise their market value, convex \omterm{def:b2:investment:adjustment-cost}{adjustment costs}, and (for the test) constant returns, so that
average $q$, observed in stock prices, equals marginal $q$.
\predicts Investment rises with $q$, and given $q$, nothing else, such as cash flow or sales, helps predict investment.
\tested Within Grunfeld's firms, holding the capital stock fixed, investment rises with market value (the fixed-effects
regression of \cref{exo:b2:investment:7}, 0.110 per dollar of value); the opening chart illustrates the same correlation.
Market value is not $q$ here, so this is consistent with the theory rather than a test of it; the test of its second
prediction is in the next field.
\breaks The fit is low, and cash flow predicts investment given $q$, especially for some groups of firms (Fazzari,
Hubbard and Petersen 1988); whether that reflects finance or mismeasured $q$ is debated (see the box below).
\end{model}

\section{The \texorpdfstring{$q$}{q} phase diagram}\label{sec:b2:investment:phase}

Two equations describe the dynamics. The capital stock grows when $q > 1$: $\dot K = K(q - 1)/\chi$ by \eqref{eq:b2:investment:q}. The value of installed capital earns its
marginal product plus capital gains, and must yield the required return: $(r + \delta)q = \mathit{MPK}(K) + \dot q$, near the steady state, where the
adjustment-cost term is negligible. So $\dot q = 0$ where $q = \mathit{MPK}(K)/(r + \delta)$, a downward-sloping curve, and $\dot K = 0$ where $q = 1$. In the steady state $\mathit{MPK}(K^*) = r + \delta$: the \omterm{def:b2:investment:desired-capital}{desired
capital stock} of the user-cost section.

\begin{omfigure}
\begin{tikzpicture}[econ, scale=0.95]
  \draw[->] (0,0) -- (7,0) node[below] {$K$};
  \draw[->] (0,0) -- (0,4.6) node[left] {$q$};
  \draw[omThm, thick] (0,2) -- (6.6,2) node[right, font=\footnotesize] {$\dot K = 0$};
  \draw[omDef, thick, domain=1:6.4, samples=40] plot (\x, {6/(\x+0.5)}) node[right, font=\footnotesize] {$\dot q = 0$};
  \draw[omMeth, very thick, domain=1.1:6.4, samples=40] plot (\x, {2 + 0.9*(2.5-\x)/(\x-0.3)});
  \node[omMeth, font=\footnotesize, anchor=west] at (1.85,3.45) {saddle path};
  \draw[omMeth, thin] (1.85,3.4) -- (1.3,3.1);
  \node[eqpoint] at (2.5,2) {};
  \draw[guide] (2.5,2) -- (2.5,0) node[below] {$K^*$};
  \node[left] at (0,2) {1};
  \draw[->, thick] (1.2,1.1) -- (1.2,0.65);
  \draw[->, thick] (1.2,1.1) -- (0.75,1.1);
  \draw[->, thick] (5.3,3.4) -- (5.3,3.85);
  \draw[->, thick] (5.3,3.4) -- (5.75,3.4);
  \draw[->, thick] (5.3,1.5) -- (5.3,1.95);
  \draw[->, thick] (5.3,1.5) -- (4.85,1.5);
  \draw[->, thick] (1.6,2.45) -- (1.6,2.0);
  \draw[->, thick] (1.6,2.45) -- (2.05,2.45);
\end{tikzpicture}

{\small The $q$ phase diagram. Above $q = 1$ the capital stock grows; above the $\dot q = 0$ curve $q$ rises. Only the saddle path, which slopes
down through the steady state, leads to it: a firm with little capital has a high $q$, invests, and $q$ falls as capital
accumulates.}
\end{omfigure}

A firm with less capital than $K^*$ has a high marginal product, so a high $q$, and invests; as capital builds up, $q$ falls back to 1. The
diagram shows what a change in policy or prices does. Take a permanent cut in the tax on profits, announced and in force
at once: the after-tax marginal product rises, and the $\dot q = 0$ curve shifts up. Capital cannot jump, but $q$, a price set in a
market, can: it jumps to the new saddle path, and the firm invests until capital reaches its new, higher steady state.

\begin{omfigure}
\begin{tikzpicture}[econ, scale=1.15]
  \draw[->] (0,0) -- (7,0) node[below] {$K$};
  \draw[->] (0,0) -- (0,4.6) node[left] {$q$};
  \draw[omThm, thick] (0,2) -- (6.6,2) node[right, font=\footnotesize] {$\dot K = 0$};
  \draw[omDef, thick, domain=1:6.4, samples=40] plot (\x, {6/(\x+0.5)});
  \node[omDef, font=\footnotesize, anchor=west] at (6.4,0.75) {before};
  \draw[omDef, thick, dashed, domain=1.6:6.4, samples=40] plot (\x, {8/(\x+0.5)});
  \node[omDef, font=\footnotesize, anchor=west] at (6.4,1.15) {after};
  \draw[omMeth, very thick, domain=2.5:5.6, samples=40] plot (\x, {2 + 0.6*(3.5-\x)/(\x-1.5)});
  \node[omMeth, font=\footnotesize, anchor=south west] at (4.3,2.0) {new saddle path};
  \node[eqpoint] at (2.5,2) {};
  \node[eqpoint] at (3.5,2) {};
  \node[font=\footnotesize, below] at (2.5,1.95) {$E_0$};
  \node[font=\footnotesize, below] at (3.6,1.95) {$E_1$};
  \node[left] at (0,2) {1};
  \draw[->, thick, black!70] (2.5,2.06) -- (2.5,2.56);
  \node[font=\footnotesize, anchor=east] at (2.3,2.3) {jump};
  \draw[->, thick, black!70] (2.62,2.66) -- (3.3,2.2);
\end{tikzpicture}

{\small A permanent cut in the tax on profits. The $\dot q = 0$ curve shifts up; the capital stock starts at $E_0$, $q$ jumps at once to the new
saddle path, and capital and $q$ move together along it to the new steady state $E_1$, with $q$ back at 1.}
\end{omfigure}

The same diagram gives the other cases. A cut announced today but in force only next year makes $q$ jump today, by less, and
investment start before the cut. A temporary cut raises $q$ by less than a permanent one, and so raises investment by less:
the firm will not install capital that costs adjustment to install and then wants to shed.

\section{Irreversible investment under uncertainty}\label{sec:b2:investment:irreversible}

Much capital, once installed, cannot be resold except at a large loss: a berth, a mine, a machine built for one plant.
Then investing today has a cost that the user cost misses: it gives up the chance of waiting to see.

\begin{definition}[Option value of waiting]\label{def:b2:investment:option-value}
The \emph{option value of waiting}\index{option value of waiting} is the amount by which the value of being able to
decide later on an irreversible investment exceeds the value of deciding now.
\end{definition}

Take a project that costs 100, sunk once spent, earns a profit $x_0$ this year, and from next year on earns $x_0(1 + u)$ or $x_0(1 - u)$ for ever,
equally likely; $r = 0.10$. Investing now is worth $11x_0 - 100$, since 1 a year for ever from today is worth $1.1/0.1 = 11$. The net present value is positive
when $x_0 > 9.09$. Waiting a year, the firm invests only in the good state. With $x_0 = 10$ and $u = 0.5$, investing now is worth 10; waiting is worth
$\tfrac12(11 \times 15 - 100)/1.1 = 29.55$, since in the bad state, with profit 5, the firm does not invest. Though the project has a positive net present value, the firm waits.
The listing finds, for each degree of uncertainty $u$, the profit above which the firm invests at once.

\omcode{code/bachelor-2/20-investment/listings/wait.py}{1}{31}{The profit that triggers an irreversible investment, for
four degrees of uncertainty.}\label{lst:b2:investment:wait}

\omoutput{code/bachelor-2/20-investment/outputs/wait.txt}{With no uncertainty the firm invests as soon as the net present
value is positive; with more uncertainty it waits for a higher profit.}\label{out:b2:investment:wait}

% code: listings/wait.py
\begin{omfigure}
\begin{tikzpicture}[econ]
\begin{axis}[omrecord, grid=none, xmin=0, xmax=18, ymin=-40, ymax=110, width=0.62\linewidth,
    height=0.42\linewidth, xlabel={Profit this year, $x_0$}, ylabel={Value}, xtick={0,9.09,15.58},
    xticklabels={0,9.09,15.58}, legend style={at={(1.03,0.5)}, anchor=west}]
  \addplot[omDef, very thick, domain=0:18, samples=2] {11*x - 100};
  \addlegendentry{Invest now}
  \addplot[omThm, very thick, domain=0:18, samples=80] {max(0, (16.5*x - 100)/2.2) + max(0, (5.5*x - 100)/2.2)};
  \addlegendentry{Wait a year}
  \draw[black!40] (axis cs:0,0) -- (axis cs:18,0);
  \node[eqpoint] at (axis cs:15.584,71.43) {};
  \draw[guide] (axis cs:15.584,71.43) -- (axis cs:15.584,-40);
  \draw[guide] (axis cs:9.09,0) -- (axis cs:9.09,-40);
\end{axis}
\end{tikzpicture}

{\small The option-value threshold, $u = 0.5$. Investing now has a positive net present value above a profit of 9.09, but waiting is worth
more until the profit reaches 15.58, where the two lines cross.}
\end{omfigure}

The threshold rises with uncertainty: 11.48 when $u = 0.25$, 15.58 when $u = 0.5$, 24.24 when $u = 0.75$. Two predictions follow. In more uncertain times firms invest
less, and they respond less to a given rise in demand, because demand must cross a higher threshold before they act.

\begin{model}{Irreversible investment}
\assumes Investment that cannot be undone without loss, uncertainty about future profits that time resolves, and the
possibility of delay.
\predicts Firms wait for profits well above the level at which the net present value turns positive; higher uncertainty
lowers investment and its response to demand.
\tested In a panel of firms, moving from the lower to the upper quartile of uncertainty typically halves the first-year
response of investment to a demand shock (Bloom, Bond and Van Reenen 2007).
\breaks When investment can be resold, when waiting lets a competitor in first, or when the uncertainty is not resolved by
waiting: then the option value shrinks.
\end{model}

\section{Financing constraints}\label{sec:b2:investment:finance}

\begin{definition}[Financing constraint]\label{def:b2:investment:financing-constraint}
A firm faces a \emph{financing constraint}\index{financing constraint} when external funds, from lenders or new shares,
cost it more than internal funds, so that its investment depends on its own cash flow as well as on the value of the
investment.
\end{definition}

If lenders know less about a project than the firm, as in \cref{ch:b2:adverse-selection}, external finance carries a premium, and a firm with
little cash may pass up projects with $q > 1$. The test regresses investment on $q$ and on cash flow: under $q$ theory alone, cash flow
should not matter.

\begin{debate}[Do financing constraints drive investment?]
% ledger: b2:fhp1988, b2:kz1997
Fazzari, Hubbard and Petersen (1988) found that investment responded strongly to cash flow, given $q$, among fast-growing firms
that paid no dividends, whose $q$ stayed high for long periods: firms, they argued, that could not raise outside funds.
Kaplan and Zingales (1997) analysed in depth the 49 low-dividend firms of that study: in only 15 per cent of
firm-years was there some question whether the firm could find funds to invest more, and the firms that appeared least
constrained had the greatest sensitivity of investment to cash flow. A high sensitivity, they concluded, cannot be read as
evidence of a constraint. The two readings differ on what the cash-flow coefficient measures; the evidence so far does not
settle how much investment \omterm{def:b2:investment:financing-constraint}{financing constraints} hold back, and for which firms.
\end{debate}

\section{Inventories and housing}\label{sec:b2:investment:housing}

Inventories, the goods firms hold in stock, are a small part of investment but a large part of its swings, since a firm
that expects sales to fall cuts production below sales and runs its stocks down. Housing is capital owned by households,
and the $q$ logic applies to it: when house prices are high relative to the cost of building, construction rises.

\begin{record}[Housing investment around 2006]
% figdata: house_summary (rise_96_06 49.4, fall_06_12 -25.8, iy_2006 22.9, iy_2010 18)
In the United States, house prices (Jord\`a--Schularick--Taylor database) deflated by the BLS consumer price index rose by 49 per cent between 1996 and 2006, then fell by 26
per cent by 2012. Total investment, of which housing is a part, went from 22.9 per cent of GDP in 2006 to 18.0 per cent in 2010.
\end{record}

The record is one episode. Over a longer run, the share of output that economies invest has varied widely.

\begin{omchart}{Gross investment as a share of GDP, per cent, United States, Germany and Japan, 1870--2020; gaps where
  the series are missing.}{Jord\`a--Schularick--Taylor Macrohistory Database (iy)}
\begin{tikzpicture}
\begin{axis}[omrecord, width=0.8\linewidth, height=0.42\linewidth, xmin=1870, xmax=2020, ymin=0, ymax=40,
    xlabel={Year}, ylabel={Per cent of GDP}, xtick={1880,1900,1920,1940,1960,1980,2000,2020},
    xticklabel style={/pgf/number format/1000 sep=}, unbounded coords=jump, legend style={at={(0.5,-0.3)},
    anchor=north, legend columns=3, /tikz/every even column/.append style={column sep=10pt}}]
  \addplot[omDef, thick] table[x=year, y=USA] {figdata/bachelor-2/20-investment/out/iy.csv};
  \addplot[omThm, thick, dashed] table[x=year, y=DEU] {figdata/bachelor-2/20-investment/out/iy.csv};
  \addplot[omMeth, thick, dotted] table[x=year, y=JPN] {figdata/bachelor-2/20-investment/out/iy.csv};
  \legend{United States, Germany, Japan}
\end{axis}
\end{tikzpicture}
\end{omchart}

% figdata: iy_summary (USA 16.9 then 20.3; DEU 12.8 then 23.4, peak 29.8 in 1964; JPN 15.5 then 27.7, peak 36.5 in 1973)
Before 1914 all three invested between about 13 and 17 per cent of GDP on average; since 1950 the averages are 20.3 per cent in the
United States, 23.4 in Germany and 27.7 in Japan, where the share peaked at 36.5 per cent in 1973.

\section{Exercises}\label{sec:b2:investment:exercises}

\begin{exercise}[$\star$]\label{exo:b2:investment:1}
Read the opening chart. What does $q$ theory predict for the relation it shows, and what in the chart is not explained by
market value?
\end{exercise}

\begin{exercise}[$\star$]\label{exo:b2:investment:2}
A machine costs 200; the interest rate is 4 per cent, depreciation 6 per cent, and machine prices rise 1 per cent a year.
Compute the user cost.
\end{exercise}

\begin{exercise}[$\star$]\label{exo:b2:investment:3}
Read the chart of investment shares. When did Japan's share peak, and how do the pre-1914 and post-1950 averages compare?
\end{exercise}

\begin{exercise}[$\star$]\label{exo:b2:investment:4}
What observation in a panel of firms would falsify the prediction of $q$ theory that, given $q$, nothing else predicts
investment? Does the debate box report one?
\end{exercise}

\begin{exercise}[$\star\star$]\label{exo:b2:investment:5}
With $\alpha = 0.3$, $Y = 1{,}000$ and user cost 0.15, the \omterm{def:b2:investment:desired-capital}{desired capital stock} is 2,000. Predict, then compute, the desired stock if the interest
rate rises by two points.
\end{exercise}

\begin{exercise}[$\star\star$]\label{exo:b2:investment:6}
With $\chi = 10$ and $\delta = 0.08$, find the investment rate when $q = 1.3$ and when $q = 0.9$. What does a $q$ below 1 mean?
\end{exercise}

\begin{exercise}[$\star\star$]\label{exo:b2:investment:7}
\textbf{(D)} On Grunfeld's data, regress investment on market value and capital, first pooled, then with a constant for
each firm. Report the coefficients and explain why they differ.
\end{exercise}

\begin{exercise}[$\star\star$]\label{exo:b2:investment:8}
Use the phase diagram to predict the path of $q$ and investment after a permanent rise in profits announced today to take
effect in one year.
\end{exercise}

\begin{exercise}[$\star\star$]\label{exo:b2:investment:9}
In the project of the irreversible-investment section, let $u = 0.25$ and $x_0 = 10$. Should the firm invest now?
\end{exercise}

\begin{exercise}[$\star\star\star$]\label{exo:b2:investment:10}
Derive \eqref{eq:b2:investment:q} from the cost of investment $I + \tfrac{\chi}{2}(I/K - \delta)^2K$ and the value $q$ of a unit of installed capital.
\end{exercise}

\begin{exercise}[$\star\star\star$]\label{exo:b2:investment:11}
For the project of the irreversible-investment section, show that while the bad state does not invest, the trigger
profit is $x^* = C(1 - 1/k)/(L(1 - (1 + u)/k))$, where $C$ is the cost of the project, $k = 2(1 + r)$ and $L = (1 + r)/r$, and that it rises with $u$.
\end{exercise}

\begin{exercise}[$\star\star\star$]\label{exo:b2:investment:12}
Suppose cash flow today signals higher profits in future years, and the stock market's $q$ measures marginal $q$ with error.
Explain why cash flow may then predict investment given measured $q$ even for a firm with no \omterm{def:b2:investment:financing-constraint}{financing constraint}.
\end{exercise}

\section{Problem: The Brenmoor container terminal}\label{sec:b2:investment:problem}

\begin{problem}[{Weekend problem --- the Brenmoor container terminal decides how many cranes to own, how fast to buy them, and
when to build an irreversible new berth}]
\label{pb:b2:investment:1}
Amounts are in millions of dalens.

\medskip
\noindent\textbf{Part I --- The user cost.} A crane costs 50; the interest rate is 6 per cent, cranes depreciate at 8 per cent a
year and crane prices rise 2 per cent a year. The marginal revenue of the $K$th crane is $20 - K$ a year.
\begin{enumerate}
  \item What is the user cost of a crane?
  \item How many cranes does the terminal want?
  \item The interest rate rises to 8 per cent. What are the user cost and the desired stock?
  \item Instead, crane prices stop rising. What is the user cost?
  \item Instead, a subsidy lowers the price of a crane to 45. What is the user cost, and the desired stock?
  \item Why does depreciation enter the user cost even though no money is paid out for it?
\end{enumerate}

\medskip
\noindent\textbf{Part II --- Adjustment costs and $q$.} The terminal owns 10 cranes; its shares are worth 900. Take $\chi = 10$ and
$\delta = 0.08$.
\begin{enumerate}[resume]
  \item What is its average $q$?
  \item Under what condition does average $q$ equal marginal $q$?
  \item Taking them equal, what is its investment rate, and how many cranes a year does it buy?
  \item What would it invest at $q = 1$?
  \item At $q = 0.8$?
  \item Why does the terminal not buy the four extra cranes it wants at once?
\end{enumerate}

\medskip
\noindent\textbf{Part III --- An irreversible berth.} A new berth costs 100, cannot be resold, and earns 12 this year. From next
year it earns 18 or 6 a year for ever, equally likely. The interest rate is 10 per cent.
\begin{enumerate}[resume]
  \item What is the net present value of building now?
  \item What is the value of waiting a year?
  \item What is the \omterm{def:b2:investment:option-value}{option value of waiting}?
  \item What should the terminal do?
  \item If the berth could be resold at its cost, what would it do?
  \item If next year's earnings were 15 or 9 instead, what would waiting be worth, and what should it do?
\end{enumerate}

\medskip
\noindent\textbf{Part IV --- The trigger.} Earnings next year are $x_0(1 + u)$ or $x_0(1 - u)$.
\begin{enumerate}[resume]
  \item At what profit $x_0$ is the net present value of building now zero?
  \item With $u = 0.5$, at what profit does the terminal build at once?
  \item With $u = 0.75$?
  \item Why do several terminals that face a rise in uncertainty together lower aggregate investment?
  \item What evidence would contradict the prediction of Question 22?
  \item State the result: the \omterm{def:b2:investment:option-value}{option value of waiting} at a profit of 12, and the profit that triggers investment, against
    the net-present-value rule.
\end{enumerate}
\end{problem}
